% Options for packages loaded elsewhere
\PassOptionsToPackage{unicode}{hyperref}
\PassOptionsToPackage{hyphens}{url}
\documentclass[
]{article}
\usepackage{xcolor}
\usepackage{amsmath,amssymb}
\setcounter{secnumdepth}{-\maxdimen} % remove section numbering
\usepackage{iftex}
\ifPDFTeX
  \usepackage[T1]{fontenc}
  \usepackage[utf8]{inputenc}
  \usepackage{textcomp} % provide euro and other symbols
\else % if luatex or xetex
  \usepackage{unicode-math} % this also loads fontspec
  \defaultfontfeatures{Scale=MatchLowercase}
  \defaultfontfeatures[\rmfamily]{Ligatures=TeX,Scale=1}
\fi
\usepackage{lmodern}
\ifPDFTeX\else
  % xetex/luatex font selection
\fi
% Use upquote if available, for straight quotes in verbatim environments
\IfFileExists{upquote.sty}{\usepackage{upquote}}{}
\IfFileExists{microtype.sty}{% use microtype if available
  \usepackage[]{microtype}
  \UseMicrotypeSet[protrusion]{basicmath} % disable protrusion for tt fonts
}{}
\makeatletter
\@ifundefined{KOMAClassName}{% if non-KOMA class
  \IfFileExists{parskip.sty}{%
    \usepackage{parskip}
  }{% else
    \setlength{\parindent}{0pt}
    \setlength{\parskip}{6pt plus 2pt minus 1pt}}
}{% if KOMA class
  \KOMAoptions{parskip=half}}
\makeatother
\setlength{\emergencystretch}{3em} % prevent overfull lines
\providecommand{\tightlist}{%
  \setlength{\itemsep}{0pt}\setlength{\parskip}{0pt}}
\usepackage{microtype}
\usepackage{titlesec}
\usepackage{booktabs}
\usepackage{array}
\usepackage{longtable}
\usepackage{ragged2e}
\usepackage{etoolbox}
\usepackage{caption}
\usepackage{fancyhdr}
\usepackage{setspace}
\usepackage{newunicodechar}

\newunicodechar{✓}{$\checkmark$}
\newunicodechar{∈}{$\in$}
\newunicodechar{∃}{$\exists$}
\newunicodechar{≠}{$\neq$}
\newunicodechar{≤}{$\le$}
\newunicodechar{≥}{$\ge$}
\newunicodechar{→}{$\rightarrow$}
\newunicodechar{₁}{$_1$}
\newunicodechar{₂}{$_2$}
\newunicodechar{ₖ}{$_k$}
\newunicodechar{─}{-}
\newunicodechar{└}{+}
\newunicodechar{├}{+}
\newunicodechar{│}{|}

\setstretch{1.03}
\setlength{\parindent}{0pt}
\setlength{\parskip}{0.45em}
\setlength{\emergencystretch}{3em}

\titleformat{\section}{\Large\bfseries}{\thesection.}{0.6em}{}
\titleformat{\subsection}{\large\bfseries}{\thesubsection.}{0.6em}{}
\titleformat{\subsubsection}{\normalsize\bfseries}{\thesubsubsection.}{0.6em}{}

\titlespacing*{\section}{0pt}{1.4em}{0.6em}
\titlespacing*{\subsection}{0pt}{1.1em}{0.4em}
\titlespacing*{\subsubsection}{0pt}{0.8em}{0.3em}

\setlength{\LTpre}{0.35em}
\setlength{\LTpost}{0.35em}
\setlength{\LTleft}{0pt}
\setlength{\LTright}{0pt}
\setlength{\tabcolsep}{4pt}
\renewcommand{\arraystretch}{1.08}
\AtBeginEnvironment{longtable}{\small}
\newcolumntype{L}[1]{>{\RaggedRight\arraybackslash}p{#1}}
\newcolumntype{C}[1]{>{\Centering\arraybackslash}p{#1}}
\newcolumntype{R}[1]{>{\RaggedLeft\arraybackslash}p{#1}}

\captionsetup[table]{skip=4pt}
\captionsetup[figure]{skip=4pt}

\pagestyle{fancy}
\fancyhf{}
\fancyfoot[C]{\thepage}
\renewcommand{\headrulewidth}{0pt}
\renewcommand{\footrulewidth}{0pt}
\setlength{\headheight}{14pt}

\newcommand{\codeid}[1]{{\small\begingroup\urlstyle{tt}\nolinkurl{#1}\endgroup}}
\usepackage{bookmark}
\IfFileExists{xurl.sty}{\usepackage{xurl}}{} % add URL line breaks if available
\urlstyle{same}
\hypersetup{
  hidelinks,
  pdfcreator={LaTeX via pandoc}}

\author{}
\date{}

\begin{document}

\section{Coverage-Preserving Compilation of Normative and Empirical
Security
Knowledge}\label{coverage-preserving-compilation-of-normative-and-empirical-security-knowledge}

\textbf{Pedro Farinha}\\
Independent Researcher\\
\href{mailto:pedro.farinha@shiftleft.pt}{\nolinkurl{pedro.farinha@shiftleft.pt}}\\
\href{https://github.com/pedrofarinhaatshiftleftpt}{github.com/pedrofarinhaatshiftleftpt}

\emph{Supported by Shiftleft - Secure Software Engineering, lda.}

\begin{center}\rule{0.5\linewidth}{0.5pt}\end{center}

\subsection{Abstract}\label{abstract}

Application security knowledge exists in two forms that are rarely
integrated: \textbf{normative knowledge} (what frameworks prescribe) and
\textbf{empirical knowledge} (what practitioners have learned).
Frameworks such as SSDF, ASVS, and SLSA define requirements but cannot
capture implementation specifics, organizational context, or practices
that predate formalization. Practitioner documentation captures these
but lacks cross-framework traceability and coverage awareness.

We present a method for \textbf{compiling} both knowledge types into a
unified, coverage-aware model. The method uses a normalized ontology
(AppSec Core {[}1{]}) as the integration substrate: normative
requirements from external frameworks and empirical content from a
practitioner corpus are independently mapped to a shared set of control
objectives, then compared to assess coverage, identify gaps, and detect
empirical extensions (content that exists in practice but is not
required by any framework).

We apply the method to a case study: a 15-chapter practitioner-authored
security manual (4,139 structural units) compiled against 5 external
frameworks (91 mapped items). The compilation reveals: (a) 78\% direct
observed coverage, rising to 95\% of in-scope items after claim-gap
resolution --- with only 4 genuine content gaps; (b) a \textbf{claim
gap} phenomenon --- most apparent gaps are traceability failures rather
than content deficiencies; and (c) measurable \textbf{empirical delta}
--- the practitioner corpus contains content that extends beyond
framework requirements, particularly in operational mechanisms and
evidence patterns.

We contribute: (1) a dual-source knowledge compilation method with
explicit coverage semantics; (2) the formal distinction between
\emph{claim gaps}, \emph{content gaps}, and \emph{cross-reference gaps};
(3) empirical evidence that practitioner knowledge substantially covers
and extends normative frameworks; and (4) the \emph{ontology-driven
refinement loop}, where compilation outputs improve the source corpus.

\textbf{Keywords:} knowledge compilation, application security,
normative knowledge, empirical knowledge, coverage analysis, security
requirements, ontology

\begin{center}\rule{0.5\linewidth}{0.5pt}\end{center}

\subsection{1. Introduction}\label{introduction}

Organizations building application security programs draw on two
distinct knowledge sources. \textbf{Normative knowledge} --- what
external frameworks prescribe --- provides structured requirements: NIST
SSDF {[}2{]} defines 20 development practices, OWASP ASVS {[}3{]}
specifies 443 verification requirements, SLSA {[}4{]} mandates build
integrity properties. \textbf{Empirical knowledge} --- what
practitioners have learned --- provides operational depth: coding
guidelines, CI/CD security configurations, container hardening
procedures, threat modeling patterns, and training curricula accumulated
through real engagements.

These two knowledge types are complementary but rarely integrated.
Frameworks provide breadth and regulatory alignment; practitioner
documentation provides depth and operational specificity. Neither alone
is sufficient:

\begin{itemize}
\tightlist
\item
  Frameworks without empirical grounding produce compliance checklists
  disconnected from implementation reality
\item
  Practitioner knowledge without framework alignment produces
  operationally rich documentation with no coverage assurance
\end{itemize}

The challenge is to \textbf{compile} both into a unified model that
preserves the coverage properties of frameworks while retaining the
operational richness of practitioner knowledge --- and to do so with
measurable, auditable properties.

\subsubsection{The Compilation Problem}\label{the-compilation-problem}

Knowledge compilation in this context is not simple merging. It
requires:

\begin{enumerate}
\def\labelenumi{\arabic{enumi}.}
\tightlist
\item
  \textbf{Normalization}: projecting heterogeneous frameworks and
  practitioner content into a shared canonical representation where
  semantic equivalence can be evaluated
\item
  \textbf{Deduplication}: recognizing when a framework requirement and a
  practitioner practice address the same concern
\item
  \textbf{Coverage assessment}: determining which normative requirements
  are satisfied by the empirical corpus
\item
  \textbf{Gap identification}: distinguishing genuine content gaps
  (nothing exists) from visibility gaps (content exists but is not
  exposed)
\item
  \textbf{Delta detection}: identifying empirical content that extends
  beyond what any framework requires
\end{enumerate}

This paper presents a method that addresses all five, using a normalized
ontology (AppSec Core {[}1{]}) as the integration substrate.

\subsubsection{Contributions}\label{contributions}

\begin{enumerate}
\def\labelenumi{\arabic{enumi}.}
\item
  \textbf{A dual-source compilation method} that integrates normative
  and empirical security knowledge through a shared ontological layer,
  with explicit coverage semantics.
\item
  \textbf{The formal distinction between claim gaps, content gaps, and
  cross-reference gaps} --- three categories of apparent non-coverage
  with different remediation strategies.
\item
  \textbf{Empirical evidence} from a case study (15-chapter manual,
  4,139 units, 5 frameworks, 91 mapped items) showing that practitioner
  knowledge substantially covers and extends normative frameworks.
\item
  \textbf{The ontology-driven refinement loop}: compilation outputs (gap
  classifications, source references) feed back into the empirical
  corpus, improving its traceability surface.
\end{enumerate}

\begin{center}\rule{0.5\linewidth}{0.5pt}\end{center}

\subsection{2. Related Work}\label{related-work}

\subsubsection{2.1 Security Framework
Alignment}\label{security-framework-alignment}

Organizations typically align practices with frameworks through manual
mapping: reading each framework's requirements and determining whether
internal documentation addresses them. Several tools and methodologies
support this (e.g., NIST SP 800-53 overlays {[}5{]}, compliance
matrices). These operate on a per-framework basis with no shared
semantic model across frameworks --- the quadratic mapping problem
identified in {[}1{]}.

\subsubsection{2.2 Knowledge Compilation in
SE}\label{knowledge-compilation-in-se}

Knowledge compilation --- transforming dispersed knowledge into a
unified, queryable representation --- has been studied in software
engineering through experience factories {[}6{]} and knowledge
management frameworks. Basili et al.~{[}6{]} established the paradigm of
accumulating and packaging practitioner experience. Our work extends
this by adding \emph{coverage-aware} compilation: the unified
representation not only stores knowledge but enables coverage assessment
against external reference sets.

\subsubsection{2.3 Requirements
Traceability}\label{requirements-traceability}

Ramesh and Jarke {[}8{]} established reference models for requirements
traceability, identifying pre-requirements, requirements, and design as
distinct traceability layers. Goknil et al.~{[}7{]} further formalized
trace relation semantics for requirements models. Our work introduces an
additional dimension to this literature: \textbf{epistemic weight} ---
different document types carry different authority, and coverage
analysis must account for this stratification. A normative requirement
(strong weight) and an operational addon (medium weight) both contribute
to coverage, but differently.

\subsubsection{2.4 Gap Analysis in
Security}\label{gap-analysis-in-security}

Gap analysis in security compliance typically compares an organization's
documented controls against a framework's requirements. Without a
structured corpus index, a common working assumption is that absence of
documentation implies absence of practice. Empirical studies of
traceability practice show that this assumption frequently fails:
practitioners omit or defer explicit traceability links for reasons
unrelated to content absence {[}9{]}. We show a specific instantiation
of this failure in the security compliance context: content frequently
exists but is not visible in the traceability surface --- a phenomenon
we formalize as \emph{claim gaps}.

\begin{center}\rule{0.5\linewidth}{0.5pt}\end{center}

\subsection{3. Dual-Source Knowledge
Model}\label{dual-source-knowledge-model}

\subsubsection{3.1 Normative Knowledge}\label{normative-knowledge}

Normative knowledge originates from external frameworks, standards, and
regulations. In our model, it consists of \textbf{external requirements}
--- prescriptive statements that define what must be done, verified, or
evidenced. Each requirement has:

\begin{itemize}
\tightlist
\item
  An \textbf{identifier} within its source framework (e.g., SSDF PW.5,
  ASVS \texttt{injection\_and\_sanitization})
\item
  A \textbf{scope} (what security concern it addresses)
\item
  A \textbf{granularity level} (practice-level, requirement-level, or
  control-level)
\end{itemize}

Normative knowledge is authoritative but incomplete: it specifies
\emph{what} but rarely \emph{how}, and it cannot capture practices that
predate the framework's publication.

\subsubsection{3.2 Empirical Knowledge}\label{empirical-knowledge}

Empirical knowledge originates from practitioner documentation: security
manuals, coding guidelines, pipeline configurations, deployment
procedures, training materials. In our model, the empirical corpus is
decomposed into \textbf{structural units} --- section-level segments,
each annotated with:

\begin{itemize}
\tightlist
\item
  \texttt{document\_role}: the type of document (requirements catalog,
  operational addon, lifecycle user story, etc.)
\item
  \texttt{normative\_weight}: the epistemic authority (strong, medium,
  low)
\item
  \texttt{heading\_path}: the hierarchical position within the corpus
\end{itemize}

This annotation enables \textbf{stratified analysis}: coverage
assessment can query the full corpus while distinguishing normative
requirements (strong weight) from operational guidance (medium weight)
from structural content (low weight/high authority).

\subsubsection{3.3 The Integration
Substrate}\label{the-integration-substrate}

Both knowledge types are mapped to a shared ontological layer: AppSec
Core v0 {[}1{]}. AppSec Core v0 provides 10 domain slices with 234 typed
instances (ControlObjective, Practice, Mechanism, Artifact). The
normalization works as follows:

\begin{itemize}
\tightlist
\item
  \textbf{Normative requirements} map to AppSec Core control objectives
  (what must be achieved)
\item
  \textbf{Empirical content} maps to AppSec Core objectives, practices,
  mechanisms, and artifacts (the full type spectrum)
\end{itemize}

The richer mapping of empirical content is expected: practitioners
document not only \emph{what} they do (objectives) but \emph{how} they
do it (practices), \emph{with what} (mechanisms), and \emph{how they
prove it} (artifacts). Frameworks typically prescribe only the first.

AppSec Core is not merely a mapping target --- it is the
\textbf{semantic comparison space}. Without a shared canonical layer,
the normative mapping and the empirical mapping would not be
commensurable: one would be expressed in SSDF/ASVS vocabulary and the
other in the corpus's internal terminology. AppSec Core makes them
comparable by reducing both to the same set of typed objectives,
enabling the coverage analysis in Stage 4.

\begin{center}\rule{0.5\linewidth}{0.5pt}\end{center}

\subsection{4. Compilation Method}\label{compilation-method}

\subsubsection{4.1 Overview}\label{overview}

The compilation proceeds in five stages:

\begin{verbatim}
Stage 1: Normative indexing     — register framework requirements
Stage 2: Empirical indexing     — decompose practitioner corpus into structural units
Stage 3: Dual mapping           — map both to AppSec Core independently
Stage 4: Coverage analysis      — compare normative surface against empirical projection
Stage 5: Editorial feedback     — feed results back into the corpus
\end{verbatim}

\subsubsection{4.2 Stage 1: Normative
Indexing}\label{stage-1-normative-indexing}

Each external framework undergoes a formal \textbf{pilot} --- a
systematic registration of its requirements into the compilation model.
For each requirement, the pilot records: source framework, identifier,
description, scope, and the AppSec Core slice(s) and objective(s) it
maps to.

The pilot follows the mapping protocol defined in {[}1{]}: each
requirement is mapped with a strength indicator (\texttt{primary} or
\texttt{secondary}) and a human-authored rationale.

\subsubsection{4.3 Stage 2: Empirical
Indexing}\label{stage-2-empirical-indexing}

The practitioner corpus is decomposed into structural units using
automated indexing. In our case study, 4,139 units were extracted from a
15-chapter manual, each annotated with \texttt{document\_role},
\texttt{normative\_weight}, and \texttt{heading\_path}. The distribution
of document roles provides the empirical corpus structure:

\textbf{Table 1.} Distribution of document roles in the SbD-ToE
empirical corpus (4,139 structural units).

\begin{longtable}{@{}L{4.2cm}R{1.6cm}R{1.8cm}L{2cm}@{}}
\toprule
\textbf{Document role} & \textbf{Units} & \textbf{Share} & \textbf{Weight} \\
\midrule
\endfirsthead
\toprule
\textbf{Document role} & \textbf{Units} & \textbf{Share} & \textbf{Weight} \\
\midrule
\endhead
\bottomrule
\endfoot
\bottomrule
\endlastfoot
addon & 1,441 & 34.8\% & medium \\
instantiated\_policy & 660 & 15.9\% & medium \\
supporting\_reference & 483 & 11.7\% & low \\
legacy\_canon & 321 & 7.8\% & authority \\
aplicacao\_lifecycle & 316 & 7.6\% & strong \\
intro & 267 & 6.5\% & low \\
kpi\_catalog & 186 & 4.5\% & medium \\
maturity & 156 & 3.8\% & medium \\
advanced & 130 & 3.1\% & medium \\
policy\_reference & 109 & 2.6\% & medium \\
requirements\_catalog & 68 & 1.6\% & strong \\
Other & 2 & <0.1\% & --- \\
\textbf{Total} & \textbf{4,139} &  &  \\
\end{longtable}

The empirical corpus is substantially larger and more diverse than the
normative surface. Strong-weight units (requirements\_catalog +
aplicacao\_lifecycle) constitute only 9.3\% of the corpus; the remaining
90.7\% is operational and contextual content --- the ``long tail'' of
empirical knowledge that frameworks do not capture.

\subsubsection{4.4 Stage 3: Dual Mapping}\label{stage-3-dual-mapping}

Both normative requirements and empirical units are independently mapped
to AppSec Core:

\begin{itemize}
\tightlist
\item
  \textbf{Normative → Core}: each framework requirement maps to one or
  more control objectives (the normative surface)
\item
  \textbf{Empirical → Core}: each chapter of the practitioner corpus
  maps to one or more slices (primary + secondary); individual
  structural units are projected to objectives, practices, mechanisms,
  and artifacts (the empirical projection)
\end{itemize}

The two mappings are independent --- neither influences the other. This
independence is important: it prevents the empirical corpus from being
shaped to fit frameworks, and vice versa.

\subsubsection{4.5 Stage 4: Coverage
Analysis}\label{stage-4-coverage-analysis}

For each normative requirement \emph{r}, the compilation assesses
whether the empirical projection contains a structural unit \emph{u}
that satisfies three criteria:

\begin{quote}
\textbf{Definition 1 (Coverage).} A normative requirement \emph{r} is
\textbf{covered} by the empirical corpus \emph{C} if there exists a unit
\emph{u ∈ C} such that: (i) \emph{u.content} addresses the same security
concern as \emph{r} (scope overlap); (ii) \emph{u} provides prescriptive
or operational guidance, not merely a mention (actionability); (iii)
\emph{u.normative\_weight} is strong or medium (authority).
\end{quote}

When coverage assessment reveals a gap, it is classified:

\begin{quote}
\textbf{Definition 2 (Content Gap).} A normative requirement \emph{r}
has a \textbf{content gap} if no unit \emph{u ∈ C} satisfies
\emph{covers(u, r)}. The concern is genuinely absent from the empirical
corpus.
\end{quote}

\begin{quote}
\textbf{Definition 3 (Claim Gap).} A normative requirement \emph{r} has
a \textbf{claim gap} if \emph{∃ u ∈ C} with \emph{covers(u, r) = true},
but the corpus's traceability surface \emph{T} does not expose this
coverage: \emph{¬∃ t ∈ T} such that \emph{maps(t, r)}. The content
exists but is not visible.
\end{quote}

\begin{quote}
\textbf{Definition 4 (Cross-reference Gap).} A normative requirement
\emph{r} has a \textbf{cross-reference gap} if a covering unit exists
but in a different chapter than expected: \emph{u.bundle\_id ≠
b\_expected}.
\end{quote}

The distinction between these gap types has practical consequences:
content gaps require authoring; claim gaps require traceability repair;
cross-reference gaps require navigation links. Conflating them leads to
unnecessary authoring effort.

Critically, coverage assessment (Definition 1) is performed in the
\textbf{normalized objective space} defined by AppSec Core {[}1{]}, not
directly between framework text and empirical text. The ``same security
concern'' in criterion (i) is mediated by shared ControlObjectives: a
framework requirement and an empirical unit both map to the same
objective, and it is the objective-level match that establishes
coverage. This ensures that coverage evaluation is independent of
framework-specific vocabulary.

\subsubsection{4.6 Stage 5: Editorial
Feedback}\label{stage-5-editorial-feedback}

The compilation's gap classifications feed back into the empirical
corpus:

\begin{itemize}
\tightlist
\item
  \textbf{Claim gaps} → traceability table corrections (add source
  reference)
\item
  \textbf{Cross-reference gaps} → inter-chapter links
\item
  \textbf{Content gaps} → authoring candidates
\end{itemize}

This feedback produces an improved version of the corpus (\emph{C → C'})
with better traceability. Subsequent compilation operates on the
improved surface, reducing false gap reports. We term this the
\textbf{ontology-driven refinement loop}.

\subsubsection{4.7 Hard Rules}\label{hard-rules}

The compilation enforces three invariants:

\begin{enumerate}
\def\labelenumi{\arabic{enumi}.}
\tightlist
\item
  \texttt{do\_not\_treat\_canon\_as\_ground\_truth} --- the traceability
  tables in the corpus are editorial claims, not verified facts;
  verification is against the structural unit index
\item
  \texttt{do\_not\_write\_directly\_from\_framework\_into\_manual} ---
  framework requirements are normalized through AppSec Core, not copied
  into the corpus
\item
  \texttt{analysis\_outputs\_are\_editorial\_candidates\_not\_automatic\_updates}
  --- all changes require human review
\end{enumerate}

\subsubsection{4.8 Coverage Preservation
Property}\label{coverage-preservation-property}

We define the property that gives the method its name:

\begin{quote}
Let \emph{r} be a normative requirement, and let \emph{normalize(r)} =
\{\emph{o₁, o₂, \ldots, oₖ}\} be the set of AppSec Core
ControlObjectives to which \emph{r} maps (at the granularity and mapping
semantics of the source framework).

The compilation is \textbf{coverage-preserving} if the coverage status
assigned to \emph{r} is determined by evaluating its normalized
objective projection --- that is, by assessing whether the empirical
corpus satisfies the ControlObjectives in \emph{normalize(r)} --- rather
than by direct comparison between framework wording and empirical text.
\end{quote}

In practice, coverage is assessed through mapped ControlObjectives
according to the granularity of the source framework: SSDF maps at
practice level, ASVS at thematic cluster level, SLSA at requirement
level. The method does not require strict satisfaction of every mapped
objective; it requires that the coverage \emph{determination} operates
in the normalized space.

Coverage preservation thus depends on the correctness of the
normalization ontology (AppSec Core {[}1{]}): if the mapping from
\emph{r} to \emph{normalize(r)} is wrong, coverage assessment is wrong
regardless of the empirical content. This dependency is a strength (it
makes evaluation principled and auditable) and a limitation (it makes
the method only as good as the normalization layer).

\begin{center}\rule{0.5\linewidth}{0.5pt}\end{center}

\subsection{5. Case Study}\label{case-study}

\subsubsection{5.1 The Empirical Corpus}\label{the-empirical-corpus}

The case study uses the \textbf{SbD-ToE} (Security by Design --- Theory
of Everything) manual: a practitioner-authored security reference
covering the full SDLC in 15 chapters. The manual was developed over 10+
years of hands-on AppSec engineering across multiple organizations. The
manual was authored in Portuguese; field names, document identifiers,
and heading paths throughout this paper reflect the corpus's original
naming.

Each chapter follows a canonical structure with four document types:

\textbf{Table 2.} SbD-ToE corpus chapter structure: four document types
per chapter.

\begin{longtable}{@{}L{3.55cm}L{7.2cm}L{1.8cm}@{}}
\toprule
\textbf{Type} & \textbf{Role} & \textbf{Weight} \\
\midrule
\endfirsthead
\toprule
\textbf{Type} & \textbf{Role} & \textbf{Weight} \\
\midrule
\endhead
\bottomrule
\endfoot
\bottomrule
\endlastfoot
\texttt{requirements\_catalog} & Typed requirement families (VAL, ERR, CFG, AUT, ACC, LOG, \ldots) & Strong \\
\texttt{aplicacao-lifecycle.md} & User stories operationalizing the chapter & Strong \\
\texttt{addon/} & Technical detail, operational guidance, examples & Medium \\
\texttt{canon/} & Normative structure, traceability tables & Authority \\
\end{longtable}

The corpus was decomposed into 4,139 structural units (Section 4.3).

\subsubsection{5.2 The Normative Surface}\label{the-normative-surface}

Five external frameworks were compiled against the corpus:

\textbf{Table 3.} External frameworks compiled against the SbD-ToE
corpus.

\begin{longtable}{@{}L{4.6cm}R{1.4cm}L{4.4cm}@{}}
\toprule
\textbf{Framework} & \textbf{Items} & \textbf{Granularity} \\
\midrule
\endfirsthead
\toprule
\textbf{Framework} & \textbf{Items} & \textbf{Granularity} \\
\midrule
\endhead
\bottomrule
\endfoot
\bottomrule
\endlastfoot
NIST SSDF v1.1 & 20 & Practices \\
OWASP ASVS v5.0 & 36 & Thematic clusters \\
SLSA Build Track v1.0 & 14 & Requirements \\
CIS Controls v8.1 (AppSec scope) & 8 & Controls \\
CAPEC v3.9 (View 683) & 13 & Attack patterns \\
\textbf{Total} & \textbf{91} &  \\
\end{longtable}

\subsubsection{5.3 Methodology: Three
Passes}\label{methodology-three-passes}

The compilation was conducted in three passes of increasing rigor:

\textbf{Pass 1 --- Editorial baseline.} The corpus's traceability tables
(\texttt{25-rastreabilidade.md} in each chapter) were read as proxies
for content. This identified 36 apparent gaps across all frameworks.

\textbf{Pass 2 --- Direct reading.} The actual addon files for the four
most-pressured chapters (Security Requirements, Secure Development, IaC,
Secure Deploy) were read directly. Many apparent gaps had substantive
content in addons not cited in the traceability tables. This motivated
the claim gap / content gap distinction.

\textbf{Pass 3 --- Systematic index verification.} For each remaining
gap, the structural unit index was queried programmatically: filter by
chapter, document role, and content; record the matching source as a
\texttt{(document\_role,\ normative\_weight,\ heading\_path)} triple.

An LLM (Claude Sonnet, Anthropic, 2025) was used as an \textbf{assisted
query interface} --- formulating filter queries and surfacing candidate
units from 4,139 entries. All coverage decisions were made by human
inspection against Definition 1. The structural unit index was the
source of truth. The LLM was not used for any text generation in this
paper.

\begin{center}\rule{0.5\linewidth}{0.5pt}\end{center}

\subsection{6. Results}\label{results}

\subsubsection{6.1 Gap Classification}\label{gap-classification}

After three passes, the 36 initial apparent gaps were reclassified:

\textbf{Table 4.} Gap classification after compilation.

\begin{longtable}{@{}L{3.35cm}R{1.5cm}R{1.8cm}L{5.15cm}@{}}
\toprule
\textbf{Gap type} & \textbf{Count} & \textbf{Share (approx.)} & \textbf{Remediation} \\
\midrule
\endfirsthead
\toprule
\textbf{Gap type} & \textbf{Count} & \textbf{Share (approx.)} & \textbf{Remediation} \\
\midrule
\endhead
\bottomrule
\endfoot
\bottomrule
\endlastfoot
Claim gap & 22 & 61\% & Traceability repair \\
Cross-reference gap & 7 & 19\% & Navigation links \\
Content gap & 4 & 11\% & Authoring needed \\
Scope exclusion & 3 & 8\% & None (deliberate) \\
\end{longtable}

\textbf{Central finding}: most apparent gaps were not content
deficiencies but failures of the traceability surface to represent
content that already existed in the corpus. The compilation method
detected this because it operates on the full structural unit index
(4,139 units), not on the traceability surface alone (approximately 15
traceability tables).

\subsubsection{6.2 Framework Coverage}\label{framework-coverage}

\textbf{Table 5.} Per-framework coverage after three-pass compilation.

\begin{longtable}{@{}L{3.8cm}R{1.2cm}R{1.6cm}R{1.6cm}R{1.8cm}R{1.8cm}@{}}
\toprule
\textbf{Framework} & \textbf{Items} & \textbf{Covered} & \textbf{Claim gap} & \textbf{Content gap} & \textbf{Scope excl.} \\
\midrule
\endfirsthead
\toprule
\textbf{Framework} & \textbf{Items} & \textbf{Covered} & \textbf{Claim gap} & \textbf{Content gap} & \textbf{Scope excl.} \\
\midrule
\endhead
\bottomrule
\endfoot
\bottomrule
\endlastfoot
SSDF v1.1 & 20 & 16 & 2 & 1 & 1 \\
ASVS v5.0 (clusters) & 36 & 28 & 6 & 2 & 0 \\
SLSA v1.0 & 14 & 11 & 2 & 0 & 1 \\
CIS v8.1 (AppSec) & 8 & 6 & 1 & 0 & 1 \\
CAPEC v3.9 & 13 & 10 & 2 & 1 & 0 \\
\textbf{Total} & \textbf{91} & \textbf{71 (78\%)} & \textbf{13 (14\%)} & \textbf{4 (4\%)} & \textbf{3 (3\%)} \\
\end{longtable}

After resolving claim gaps (traceability repairs), effective coverage
rises to \textbf{84 of 88 in-scope items (95\%)} --- within this case
study, under the author-evaluated coverage protocol described in §5.3.

\subsubsection{6.3 Content Gaps (Genuine)}\label{content-gaps-genuine}

Only four genuine content gaps were confirmed:

\begin{enumerate}
\def\labelenumi{\arabic{enumi}.}
\tightlist
\item
  \textbf{SLSA L3 hardened builds} --- the corpus covers L1/L2 but lacks
  hermetic/reproducible build specifics
\item
  \textbf{CAPEC-691 (dependency metadata spoofing)} --- dependency
  confusion addressed; metadata spoofing not
\item
  \textbf{CAPEC-443 (insider developer threat)} --- branch protection
  present; authorized-insider threat not distinct
\item
  \textbf{SSDF PW.6 (compilation/interpreter security)} --- partially
  covered; compiler flags and interpreter hardening lack dedicated
  treatment
\end{enumerate}

These represent genuine authoring needs not resolvable by traceability
improvements.

\subsubsection{6.4 Representative Claim Gap
Reclassifications}\label{representative-claim-gap-reclassifications}

\textbf{SSDF PO.2 --- Roles \& Responsibilities.} Appeared as a gap in
three chapters. Content found in
\texttt{00-fundamentos/roles-responsabilidades/} --- 13 security roles
with per-chapter responsibilities. Invisible because the chapter had no
traceability table. Resolution: create traceability entry (no content
authoring).

\textbf{ASVS \texttt{secure\_configuration\_baseline}.} Appeared as a
gap in four chapters. Content found in: Cap. 02 \texttt{addon/07}
(CFG-001→007: debug off, env separation, vault, drift); Cap. 08 IaC
principles addon (OPA policies, privilege minimum); Cap. 09 OPA/Kyverno
admission addon. Three independently authored addons, never
cross-referenced.

\textbf{SSDF PW.7 --- Code Review.} Classified as ``partial --- not read
in full, assumed partial.'' Verified as covered:
\texttt{aplicacao\_lifecycle\ (strong):\ US-02\ Revisão\ de\ Código\ Segura}
--- structured review with criteria and traceability.

\subsubsection{6.5 Empirical Delta}\label{empirical-delta}

The compilation reveals empirical content that \textbf{extends beyond}
what the five frameworks require:

\begin{longtable}{@{}L{3.35cm}L{4.85cm}L{4.8cm}@{}}
\toprule
\textbf{Category} & \textbf{Example} & \textbf{Framework status} \\
\midrule
\endfirsthead
\toprule
\textbf{Category} & \textbf{Example} & \textbf{Framework status} \\
\midrule
\endhead
\bottomrule
\endfoot
\bottomrule
\endlastfoot
Operational mechanisms & OPA/Kyverno admission policies, securityContext enforcement & Not prescribed by SSDF or ASVS \\
Evidence patterns & Requirement-linked test suites, validation coverage reports & Not required at this granularity \\
Role-based training curricula & Per-chapter training paths with knowledge validation & CIS-14 partially addresses; SSDF does not \\
Proportional application & Risk-level-dependent requirement activation (L1/L2/L3) & No framework specifies proportionality \\
Exception governance & Formal exception process with lifecycle, audit trail & Not prescribed at this operational level \\
\end{longtable}

This delta is the empirical knowledge that practitioners have
accumulated but that normative frameworks have not (yet) formalized. It
represents the value that empirical knowledge adds beyond framework
compliance.

\subsubsection{6.6 Refinement Loop
Results}\label{refinement-loop-results}

The compilation directly produced editorial corrections to the corpus:

\begin{longtable}{@{}L{5.2cm}R{1.2cm}L{5.85cm}@{}}
\toprule
\textbf{Correction type} & \textbf{Count} & \textbf{Effect} \\
\midrule
\endfirsthead
\toprule
\textbf{Correction type} & \textbf{Count} & \textbf{Effect} \\
\midrule
\endhead
\bottomrule
\endfoot
\bottomrule
\endlastfoot
Traceability table status reclassifications & 22 & Partial -> semantic \\
Source reference additions (\texttt{Fonte verificada} column) & 29 & Each cites a specific \texttt{(role, weight, path)} triple \\
Cross-chapter navigation links & 7 & Content now reachable from expected chapter \\
New traceability table created & 1 & Cap. 00 had no \texttt{canon/} directory \\
\end{longtable}

The corpus improved not by writing new security content but by
\textbf{using the compilation to make existing content visible}. This is
the refinement loop in action.

\begin{center}\rule{0.5\linewidth}{0.5pt}\end{center}

\subsection{7. Discussion}\label{discussion}

\subsubsection{7.1 Why Empirical Knowledge
Matters}\label{why-empirical-knowledge-matters}

The case study demonstrates that a mature practitioner corpus covers
95\% of the normative surface (after claim gap resolution) while
containing substantial content not captured by any framework. This
suggests that:

\begin{itemize}
\tightlist
\item
  Frameworks \textbf{lag} practice: practitioners encounter and document
  security concerns before frameworks formalize them (e.g., CFG-001→007
  documented in the internal corpus in 2023, based on corpus commit
  history; ASVS \texttt{secure\_configuration\_baseline} formalized in
  2025)
\item
  Frameworks \textbf{abstract} practice: they specify \emph{what} but
  not \emph{how}, leaving the operational layer to practitioners
\item
  Empirical knowledge is not a luxury --- it is the
  \textbf{implementation substance} that frameworks assume but do not
  provide
\end{itemize}

A compilation method that captures only normative knowledge loses this
substance. A method that captures only empirical knowledge lacks
coverage assurance. The dual-source model captures both.

\subsubsection{7.2 The Claim Gap
Phenomenon}\label{the-claim-gap-phenomenon}

The predominance of claim gaps (22 of 36 apparent gaps) has a structural
cause. The corpus's traceability surface (\texttt{canon/} directory)
systematically references strong-weight units (requirements\_catalog,
aplicacao\_lifecycle) but under-references medium-weight units (addons).
Since addons constitute 34.8\% of the corpus and contain the richest
operational content, any analysis that reads only the traceability
surface misses them.

This is not a flaw of this specific corpus --- it is a structural
consequence of \textbf{epistemic stratification}: when documentation is
organized by authority level, the traceability layer tends to index the
normative layer while the operational layer remains invisible. We
conjecture that the claim gap phenomenon generalizes to other similarly
stratified corpora, given the structural cause identified --- but
cross-organizational validation is needed to test this claim (see §8).

\subsubsection{7.3 Coverage Adequacy, Not
Completeness}\label{coverage-adequacy-not-completeness}

The method assesses \textbf{coverage adequacy} --- whether the empirical
corpus addresses specific normative requirements --- not \textbf{domain
completeness} --- whether the corpus captures all possible AppSec
concerns. These are different properties:

\begin{itemize}
\tightlist
\item
  Coverage adequacy is measurable against a reference set (the
  frameworks)
\item
  Domain completeness is not finitely verifiable
\end{itemize}

Our results (95\% framework coverage, 4 content gaps) speak to the
former. The latter would require validation against a broader set of
practitioners and domains.

\subsubsection{7.4 Generalizability}\label{generalizability}

The compilation method is general: it requires (a) a normalization
ontology, (b) a decomposed empirical corpus with role/weight
annotations, and (c) an indexed normative surface. The specific ontology
(AppSec Core), corpus (SbD-ToE), and frameworks (SSDF, ASVS, etc.) are
instances. The method transfers to any domain with heterogeneous
normative sources and practitioner documentation --- though we have
validated it only in the AppSec domain.

\begin{center}\rule{0.5\linewidth}{0.5pt}\end{center}

\subsection{8. Limitations}\label{limitations}

\textbf{Single empirical corpus.} The case study uses one organization's
manual. Coverage rates and claim gap prevalence may differ for other
organizations. We argue for structural transferability of the claim gap
phenomenon but have not validated it cross-organizationally.

\textbf{Author-evaluated coverage.} The \texttt{covers(u,\ r)} function
was evaluated by the paper's authors. Inter-rater reliability was not
measured. We partially mitigate this by releasing a curated artifact set
containing the first-wave pilot manifests, normalization outputs, and
terminal comparison artifacts that support the reported case-study
surface. The full internal structural unit index and all coverage
decisions are not part of the public \texttt{v1.0.0} package.

\textbf{Practice/cluster-level granularity.} Coverage assessment
operates at SSDF practice level and ASVS cluster level, not at
individual requirement level. Finer-grained assessment would increase
precision but also evaluation cost.

\textbf{Temporal snapshot.} Framework versions are captured as of early
2026. New versions require re-compilation.

\textbf{Empirical delta not independently validated.} The delta analysis
(Section 6.5) identifies content extending beyond frameworks but does
not verify whether this content is valuable or correct --- only that it
exists.

\begin{center}\rule{0.5\linewidth}{0.5pt}\end{center}

\subsection{9. Conclusion}\label{conclusion}

We have presented a method for compiling normative and empirical
security knowledge into a unified, coverage-aware model. The method uses
a normalized ontology as the integration substrate, enabling independent
mapping of framework requirements and practitioner content to a shared
semantic space.

The case study reveals three findings:

\begin{enumerate}
\def\labelenumi{\arabic{enumi}.}
\item
  \textbf{Claim gaps dominate apparent non-coverage.} 22 of 36 apparent
  gaps were traceability failures rather than content deficiencies.
  Compilation against the full indexed corpus --- rather than the
  traceability surface alone --- is necessary to distinguish the two.
\item
  \textbf{Empirical knowledge substantially covers normative
  requirements.} After claim gap resolution, 95\% of in-scope framework
  requirements were covered, with only 4 genuine content gaps across 91
  mapped items.
\item
  \textbf{Empirical knowledge extends beyond frameworks.} The
  practitioner corpus contains operational mechanisms, evidence
  patterns, proportional application models, and governance processes
  that no analyzed framework prescribes --- the empirical delta.
\end{enumerate}

The compilation method produces a measurable, auditable integration of
both knowledge types. The ontology-driven refinement loop ensures that
compilation outputs improve the source corpus, creating a virtuous
cycle: the more the corpus is compiled against frameworks, the better
its traceability becomes.

The method is coverage-preserving with respect to the normalized
objective model: both knowledge sources --- normative and empirical ---
are projected into the same canonical objective layer, and coverage is
determined by evaluating that projection rather than by direct
framework-to-text comparison. The preservation property (Section 4.8)
depends on the correctness of the normalization ontology and the
granularity of the mapping; it is only as reliable as the underlying
normalization layer. This is a principled dependency: it makes the
assumptions explicit and the evaluation auditable.

Knowledge compilation for security is not a one-time alignment exercise.
It is an ongoing process in which normative frameworks provide coverage
targets, empirical knowledge provides implementation substance, and a
shared ontological layer enables the two to be systematically compared,
integrated, and improved.

\begin{center}\rule{0.5\linewidth}{0.5pt}\end{center}

\subsection{10. Artifact Availability}\label{artifact-availability}

Curated supporting artifacts for this paper are available in the
companion public repository at
\url{https://github.com/sbd-ai-runtime/appsec-core-ontology-research}.
For this paper, the relevant materials are organized under
\texttt{papers/02-coverage-preserving-knowledge-compilation/artifacts/},
notably
\texttt{papers/02-coverage-preserving-knowledge-compilation/artifacts/pilot\_manifests/}
and
\texttt{papers/02-coverage-preserving-knowledge-compilation/artifacts/pilot\_outputs/},
which contain the released first-wave pilot manifests and comparison
outputs supporting the case-study surface. The same repository also
contains this paper's curated source under
\texttt{papers/02-coverage-preserving-knowledge-compilation/source/} and
its public PDF under
\texttt{papers/02-coverage-preserving-knowledge-compilation/pdf/}.

\begin{center}\rule{0.5\linewidth}{0.5pt}\end{center}

\subsection{References}\label{references}

{[}1{]} P. Farinha, ``AppSec Core: A normalized ontology for security
requirements across heterogeneous frameworks,'' preprint, 2026. {[}arXiv
ID: TBD --- submitted as companion preprint{]}

{[}2{]} NIST. Secure Software Development Framework (SSDF) Version 1.1.
SP 800-218. 2022.

{[}3{]} OWASP. Application Security Verification Standard (ASVS) v5.0.0.
2025.

{[}4{]} SLSA. Supply-chain Levels for Software Artifacts. Build Track
v1.0. 2023.

{[}5{]} NIST. Security and Privacy Controls for Information Systems and
Organizations. SP 800-53 Rev.~5. 2020.

{[}6{]} V. R. Basili, F. Shull, and F. Lanubile, ``Building knowledge
through families of experiments,'' \emph{IEEE Trans. Software Eng.},
vol.~25, no. 4, pp.~456--473, Jul./Aug.~1999, doi: 10.1109/32.799939.

{[}7{]} A. Goknil, I. Kurtev, K. van den Berg, and J.-W. Veldhuis,
``Semantics of trace relations in requirements models for consistency
checking and inferencing,'' \emph{Software \& Systems Modeling},
vol.~10, no. 1, pp.~31--54, 2011, doi: 10.1007/s10270-009-0142-3.

{[}8{]} B. Ramesh and M. Jarke, ``Toward reference models for
requirements traceability,'' \emph{IEEE Trans. Software Eng.}, vol.~27,
no. 1, pp.~58--93, Jan.~2001, doi: 10.1109/32.895989.

{[}9{]} P. Mäder and J. Cleland-Huang, ``Why don't we trace? A study on
the barriers to software traceability in practice,'' \emph{Requirements
Engineering}, 2023, doi: 10.1007/s00766-023-00408-9.

\end{document}
